\documentclass{article}
\usepackage[utf8]{inputenc}
\usepackage[russian]{babel}
\usepackage{amsmath}
\usepackage{graphicx}
\usepackage{amsfonts}
\usepackage{geometry}
\geometry{a4paper, margin=1in}

\title{Описание программы построения расчётной сетки}
\author{Имя автора}
\date{}

\begin{document}

Отметим, что аппроксимация якобиана в каждом из четырех способов (5.2) с точностью до множителя $\frac{\Delta \alpha \Delta \beta}{2}$ представляет площадь одного из треугольников ячейки ( $\mathrm{i}+1 / 2, \mathrm{j}+1 / 2$ ) на плоскости ху (см. рис. 5.1 ) - способ 1 соответствует треугольнику , образованному точками ( $\mathrm{i}, \mathrm{j}$ ), ( $\mathrm{i}+1, \mathrm{j})$ и ( $\mathrm{i}, \mathrm{j}+1$ ), способ 2 - треугольнику ( $\mathrm{i}, \mathrm{j}),(\mathrm{i}+1, \mathrm{j})$ и ( $\mathrm{i}+1, \mathrm{j}+1$ ), способ 3 - треугольнику ( $\mathrm{i}+1, \mathrm{j})$, $(\mathrm{i}+1, \mathrm{j}+1)$ и $(\mathrm{i}, \mathrm{j}+1)$, и наконец способ 4 - треугольнику $(\mathrm{i}, \mathrm{j}),(\mathrm{i}+1, \mathrm{j}+1)$ и ( $\mathrm{i}, \mathrm{j}+1)$,:
Рис. 5.1
Ячейка (i+1/2, j+1/2) и составляющие ее треугольники.
Следствием этого является тот факт, что минимум дискретного функционала (5.1) при выборе усреднения по формуле (5.3) может достигаться только при положительных значениях площадей указанных треугольников, то есть построенная сетка будет иметь выпуклые ячейки. Действительно, выпуклость ячейки может потеряться только при переходе площади одного из треугольников через ноль (см. рис. 5.2), но при этом одно из слагаемых в (5.3) будет возрастать до бесконечности. Это свойство, установленное в работе [12], называют барьерным свойством функционала (5.4).

Отметим, что при других способах усреднения, например, использования для производных усредненных значений

\end{document}
